\section{Paths with an escape at each endpoint}
\label{path:section}

We retain the notation of the corridor deformation. In particular,
$\pc/2\le p<\pc$, $T_0=T_p$, and $I_w$ is the interior of the
chronological loop erasure of a raw $k$-step $P$-path $w$ from $u$ to $v$.
The domain $D_w=V\setminus I_w$ depends only on this path. The next lemma
finds paths for which deleting $I_w$ leaves useful potentials accessible
from both endpoints. All probabilities in the lemma concern auxiliary
walks, independently of the percolation configuration.

Rooted reversibility and mass transport provide the general framework
for the construction; see \cite[Theorem~4.1]{AldousLyons2007}.
The labeled, type-conditioned experiment below supplies the particular
reversal and domination identities needed here. It makes no assertion
that chronological loop erasure commutes with path reversal.

\begin{lemma}[Two endpoint escapes]
\label{path:selection}
For every $\tau>0$ there are constants $J\in\N$, $C_1>0$, $k_0\in\N$,
and $b_*>0$, depending only on $G$ and $\tau$, with the following property.
For every $\pc/2\le p<\pc$, every $k\ge k_0$, and every $R\ge1$, a set
of paths of $\E_{\mathrm{path}}$-probability at least $b_*$ has
loop-erased length at least three and admits, from each of $u,v$, a
simple path in $D_w$ of length at most
\[
 j=J+\left\lceil C_1\log((k+1)R)\right\rceil
\]
to a vertex $z$ satisfying
\begin{equation}
 \sum_{b\in\Nb[I_w]}T_0(z,b)\le\frac{\tau}{R}.
 \label{eq:s33}
\end{equation}
The constants $J,C_1,k_0,b_*$ are fixed before $p,k,R$ are chosen.
\end{lemma}

We first record a smoothing bound which does not involve $\norm{T_0}$.
It is a quasi-transitive, off-diagonal form of the random-walk averaging
estimate known as Schramm's lemma; see \cite[Section~1.1]{Kozma2011}
and \cite[Proposition~6.4]{Hutchcroft2019}. The finite trace gives a
short proof adapted to our symmetric lazy walk.

\begin{lemma}[Smoothing a connection kernel]
\label{path:smoothing}
For all vertices $a,b$ and integers $l\ge0$,
\begin{equation}
 (P^lT_0)(a,b)\le\sqrt{c_\sigma}\,\rho^l.
 \label{eq:s34}
\end{equation}
\end{lemma}

\begin{proof}
Positive-semidefinite Cauchy--Schwarz for $T_0$ gives
\[
 \abs{(P^lT_0)(a,b)}^2
 \le (P^lT_0P^l)(a,a)T_0(b,b).
\]
Here $T_0(b,b)=1$. The invariant positive operator $P^lT_0P^l$ has
every diagonal entry bounded by $c_\sigma$ times its weighted trace.
Consequently, trace cyclicity and the positive-trace bound give
\[
 (P^lT_0P^l)(a,a)
 \le c_\sigma\tr(P^{2l}T_0)
 \le c_\sigma\norm{P^{2l}}\tr(T_0)
 =c_\sigma\rho^{2l}.
\]
This argument uses no commutation between $P$ and $T_0$.
\end{proof}

\begin{proof}[Proof of \cref{path:selection}]
Let $L_{\mathrm{type}}$ and $c_*>0$ be such that every transition
probability of the finite type chain at every time $L\ge
L_{\mathrm{type}}$ is at least $c_*$. The same $c_*$ works for all
types and will therefore remain fixed as $J$ varies.

\paragraph{The labeled experiment.}
For an integer $J$ to be chosen, take a geodesic segment
$\gamma=(\gamma_{-J},\ldots,\gamma_J)$. Such segments exist because an
infinite connected locally finite graph has infinite diameter. Write
$\mathcal S=\Gamma\gamma$ for its orbit as an ordered full segment,
$O_i=\Gamma\gamma_i$, and $\vartheta_i$ for the root weight of $O_i$.
Different indices may describe the same vertex orbit. For $x\in O_i$,
let
\[
 N_i=\#\{g\in\mathcal S:g_i=x\}.
\]
This is a positive finite number independent of the choice of $x$ in
$O_i$. Indeed an anchored segment is contained in the finite ball
$B(x,2J)$, and automorphisms give bijections between anchor sets.
The full orbit $\mathcal S$ itself need not be finite.

Fix $u\in O_0$ and put $L=k-2J\ge L_{\mathrm{type}}$. Sample $g$
uniformly among segments with $g_0=u$. Follow its positive half from
$u$ to $g_J$, then run a middle $P$-walk $\eta$ of length $L$ conditioned
to end in $O_J$. Its conditioning probability
\[
 q_L=P^L(a,O_J),\qquad a\in O_J,
\]
is constant on $O_J$ and is at least $c_*$. Given $\eta$, sample $h$
uniformly among segments with $h_J=\eta_L$, and follow
$(h_J,h_{J-1},\ldots,h_0)$ to $v=h_0$. The resulting raw path $w$ has
length $k$.

Independently, after these choices, take ordinary $P$-walks $\alpha$
and $\beta$ of a common length $m$ from $g_{-J}$ and $h_{-J}$.
Prepend the corresponding negative segment to each walk. These are the
two candidate escapes, of length $J+m$; see
\cref{fig:escape-templates}.
\begin{figure}[tbp]
\centering
\begin{tikzpicture}[x=1cm,y=1cm,font=\small,
  dot/.style={circle,fill,inner sep=1.6pt},
  raw/.style={very thick,blue!65!black,-{Stealth[length=2mm]}},
  escape/.style={thick,orange!70!black,-{Stealth[length=2mm]}}]
\node[dot,label={above left:$u=g_0$}] (u) at (0,0) {};
\node[dot,label=above:$g_J$] (a) at (1.5,1.35) {};
\node[dot,label=above:$h_J$] (b) at (6.5,1.35) {};
\node[dot,label={above right:$v=h_0$}] (v) at (8,0) {};
\node[dot,label=left:$g_{-J}$] (c) at (-1.1,-1.1) {};
\node[dot,label=right:$h_{-J}$] (d) at (9.1,-1.1) {};
\node[dot,label=below:$z$] (z) at (0,-2.65) {};
\node[dot,label=below:$z'$] (zz) at (8,-2.65) {};
\draw[raw] (u)--node[above left] {$J$} (a);
\draw[raw] (a) to[out=-10,in=170] (2.6,1.1)
  to[out=-10,in=180] (4,1.3) to[out=0,in=190] (b);
\node[align=center] at (4,2.65) {middle walk of length $L$\\
conditioned to end in the type of $g_J$};
\draw[raw] (b)--node[above right] {$J$} (v);
\draw[escape] (u)--node[left] {$J$} (c);
\draw[escape] (v)--node[right] {$J$} (d);
\draw[escape] (c) to[out=-80,in=180] (-.3,-1.9) to[out=0,in=180] (.5,-2.1)
  to[out=0,in=0] (z);
\draw[escape] (d) to[out=-100,in=0] (8.3,-1.9) to[out=180,in=0] (7.5,-2.1)
  to[out=180,in=180] (zz);
\node[align=center,text width=5cm] at (4,-1.2)
{Walk tails\\of length $m=j-J$};
\end{tikzpicture}
\caption{The anchored path experiment. The blue raw path has a
prescribed prefix, a reversible middle walk, and a prescribed suffix.
The two orange escapes use the opposite halves of the same geodesic
segments and walk tails that, conditional on $g,\eta,h$, are independent
ordinary $P$-walks before avoidance is imposed. Reversing the middle
walk and exchanging the two segments and escape tails gives the reversed
raw-path law after transport on the center orbit. Avoidance is tested
against the raw path; the escapes may intersect each other. Loop-erasing
either escape keeps it outside the corridor interior.}
\label{fig:escape-templates}
\end{figure}

Write $Q_u$ for the law of the
full labeled object $\xi=(g,\eta,h,\alpha,\beta)$, and $\mu(\zeta)$
for the product of the transition probabilities along any walk $\zeta$.
Its exact probability is
\begin{equation}
 Q_u(\xi)=
 \frac{\mu(\eta)\mu(\alpha)\mu(\beta)}{N_0N_Jq_L}.
 \label{path:labeled-weight}
\end{equation}

\paragraph{Reversal and comparison with the ordinary walk.}
The labeled law serves two purposes. Reversal lets us estimate avoidance
at one endpoint and transfer the estimate to the other. Domination by
the ordinary walk law then converts successful labeled experiments into
a positive mass of the raw paths used in the deformation.

There is a rooted reversal of this experiment:
\[
 \xi^\dagger=(h,\overleftarrow\eta,g,\beta,\alpha),
\]
rooted at $v$. Symmetry of $P$ shows that
$Q_u(\xi)=Q_v(\xi^\dagger)$ by \eqref{path:labeled-weight}; neither
$N_0=N_J$ nor equality of the center and endpoint type weights is
needed. The escape tails are swapped, rather than reversed. For an
invariant nonnegative function $F$ of the labeled object, apply
\eqref{eq:s02} to
\[
 M_F(x,y)=\sum_{\xi:x\to y}Q_x(\xi)F(\xi),
 \qquad x,y\in O_0,
\]
and set this transport to zero off $O_0\times O_0$. Invariance makes
the outgoing sum constant on $O_0$, and the common root weight
$\vartheta_0$ cancels. The reversal bijection identifies the incoming
sum with the outgoing sum for $F\circ\dagger$. Thus
\begin{equation}
 \E_{Q_u}F=\E_{Q_u}(F\circ\dagger).
 \label{path:rooted-reversal}
\end{equation}

We also need domination after forgetting all labels. For a fixed raw
path $w$, let $n_g$ and $n_h$ count the full segments compatible with
its visible first and last blocks, respectively. Its middle is fixed,
and summing all compatible full segments and all escape tails gives
\[
 Q_u(w)=\frac{\mu(\eta)}{q_L}
        \frac{n_g}{N_0}\frac{n_h}{N_J}
 \le \frac{\mu(\eta)}{c_*}.
\]
The two fractions are at most one. In particular, different negative
halves which project to the same visible half cause no additional
multiplicity. The ordinary walk probability is
$\mu_u(w)=(2d)^{-2J}\mu(\eta)$, so
\begin{equation}
 Q_u(w)\le D_J\mu_u(w),\qquad
 D_J=\frac{(2d)^{2J}}{c_*}.
 \label{path:domination}
\end{equation}

\paragraph{Avoiding the raw path.}
We bound all collisions uniformly in the escape-tail length $m$.
First consider the escape from $u$, conditional on $g$. Its deterministic
negative half and the positive half of $g$ meet only at $u$, which is
allowed. Before conditioning the middle endpoint type, the probability
that middle time $r$ visits any one of the $J+1$ negative-segment
vertices is at most $\rho^r$, and is zero for $r<J$. Similarly, at
escape-tail time $l$ the probability of visiting any one of the $J+1$
positive-segment vertices is at most $\rho^l$, and is zero for $l<J$.
For a collision of the random escape tail and the middle at times
$l,r$, their independence and symmetry of $P$ give
\[
 \sum_xP^l(g_{-J},x)P^r(g_J,x)
 =P^{l+r}(g_{-J},g_J)\le\rho^{l+r}.
\]
This probability is zero when $l+r<2J$. Conditioning the middle
endpoint costs at most $c_*^{-1}$ in these estimates. Hence all these
collisions have total probability at most
\begin{equation}
 E_J=\frac1{c_*}\left(
 2(J+1)\sum_{r\ge J}\rho^r
 +\sum_{n\ge2J}(n+1)\rho^n\right),
 \qquad E_J\longrightarrow0.
 \label{path:early-collisions}
\end{equation}
Extending the time sums to infinity is what makes this bound uniform
in $L$ and $m$.

For collisions with the last block, define the anchored-position kernel
\[
 H_{a,b}(x,y)=
 \one_{\{x\in O_a\}}
 \frac{\#\{g\in\mathcal S:g_a=x,\ g_b=y\}}{N_a}.
\]
Its row sums are one on $O_a$ and zero elsewhere. Its column sum is
constant on $O_b$, is zero elsewhere, and, by \eqref{eq:s02}, is exactly
\begin{equation}
 \sum_xH_{a,b}(x,y)=\frac{\vartheta_a}{\vartheta_b}
 \quad(y\in O_b),\qquad
 \norm{H_{a,b}}\le
 \sqrt{\frac{\vartheta_a}{\vartheta_b}}\le\sqrt{c_\sigma}.
 \label{path:position-kernel}
\end{equation}
Thus the operator bound is independent of $J$ and of all anchor counts.
Conditional on $g$, the middle endpoint distribution $\nu$ satisfies
$\norm{\nu}_2\le c_*^{-1}\rho^L$. Each suffix-position distribution
$\nu H_{J,i}$, $0\le i\le J$, therefore has norm at most
$\sqrt{c_\sigma}\,c_*^{-1}\rho^L$.

Conditional on $g$, the first escape is independent of the middle and
$h$. Its deterministic positions are point masses, and the distribution
at tail time $l$ has norm at most $\rho^l$. Cauchy--Schwarz and a union
bound now show that its collisions with the suffix have probability at
most
\begin{equation}
 F_J\rho^L,\qquad
 F_J=C\left((J+1)^2+\frac{J+1}{1-\rho}\right),
 \label{path:suffix-collisions}
\end{equation}
where $C$ depends only on $G$. The estimate includes all tail times
because $\sum_{l\ge0}\rho^l=(1-\rho)^{-1}$. The same suffix-position
bound at $i=0$ gives
\[
 Q_u(\dist(u,v)<3)
 \le \abs{B(u,2)}^{1/2}\sqrt{c_\sigma}\,c_*^{-1}\rho^L
 \le C_0\rho^L,
\]
with $C_0$ independent of $J$. By \eqref{path:rooted-reversal}, the
second escape has exactly the same avoidance-failure bound as the
first: the relevant event refers only to the raw path and its reversed
raw path.

Choose $J$ so large that $2E_J\le1/8$, and then choose
$L_0\ge L_{\mathrm{type}}$ so large that
$(2F_J+C_0)\rho^{L_0}\le1/8$. For every $L\ge L_0$ and every $m$,
with probability at least $3/4$ both candidate escapes avoid every raw
path vertex other than their respective starting vertices, and
$\dist(u,v)\ge3$. The escapes need not avoid each other. This choice
order is valid because no anchor-count factor occurs in
\eqref{path:early-collisions}.

Avoidance will keep the escape paths inside $D_w$. It remains to choose
their endpoints so that the deletion error in \cref{def:deletion} is
small. For this we use the smoothing bound before imposing avoidance.

\paragraph{Testing the potentials.}
Set $\lambda=-\log\rho>0$, choose $C_1$ with $\lambda C_1>1$, and put
$m=\lceil C_1\log((k+1)R)\rceil$. Conditional on $g,\eta,h$, and
\emph{before imposing any avoidance event}, the two tails are independent
ordinary $P$-walks. Their terminal vertices $Z_1,Z_2$ therefore satisfy,
by \cref{path:smoothing} and $\abs{\Nb[I_w]}\le(d+1)(k+1)$,
\[
 \E\left[\sum_{b\in\Nb[I_w]}T_0(Z_i,b)\,\middle|\,g,\eta,h\right]
 \le C(k+1)\rho^m,\qquad i=1,2.
\]
Here the loop-erased interior is defined whether or not its length is
at least three. Markov's inequality bounds the sum of the two potential
failure probabilities by
\begin{equation}
 2C\tau^{-1}R(k+1)\rho^m
 \le 2C\tau^{-1}((k+1)R)^{1-\lambda C_1}.
 \label{path:potential-failures}
\end{equation}
Increase a fixed $k_0\ge2J+L_0$ until the right-hand side is at most
$1/4$ for all $k\ge k_0$ and $R\ge1$. All constants remain independent
of $p$ and $R$.

Subtracting the failure probability in \eqref{path:potential-failures}
from the preceding $3/4$ success bound leaves probability at least
$1/2$ for both avoidance and both potential tests simultaneously.
Since $\dist(u,v)\ge3$, the retained corridor has length at least three.
Its interior is contained in the raw path and excludes $u,v$. Each
escape therefore lies in $D_w$, and erasing its loops gives a simple
path to the same terminal vertex of length at most $J+m$. No reversal
property of chronological loop erasure is used.

Project this successful event onto raw paths admitting these witnesses.
By \eqref{path:domination}, ordinary walk law from a root in $O_0$
assigns these paths probability at least $1/(2D_J)$. Averaging over
root types proves the lemma with
$b_*=\vartheta_0/(2D_J)>0$.
\end{proof}
